\keys{Wildfires; Computer vision; Edge AI; LoRa; Nvidia Jetson Orin Nano; Environmental monitoring.} % Palavras chave em inglês separadas por ponto e vírgula


% Não deixe linha em branco após \begin{abstract}{Title in English}
\begin{abstract}{Remote Monitoring in Conservation Units}
The persistent incidence of wildfires poses a serious socio-environmental and economic threat in Brazil, particularly in the Cerrado biome, which accounted for 68.5\% of the country's burned area in 2025, predominantly within conservation units. Although technological solutions for the remote monitoring of fire hotspots based on artificial intelligence (AI), sensing, and computer vision are available, current systems face major operational bottlenecks: memory limitations on edge devices (Edge AI), connectivity constraints in remote areas without mobile coverage, and, above all, high deployment and maintenance costs, which can exceed BRL 8 million for larger-scale applications using commercial alternatives based on monitoring towers. In response to this scenario, this study proposes a low-cost, high-efficiency architecture for the early detection of wildfires in Environmental Protection Areas. The proposed solution integrates a PTZ (Pan-Tilt-Zoom) camera to scan and provide visual coverage of the monitored area, and an NVIDIA Jetson Orin Nano edge computer to run computer vision models locally, powered by an autonomous photovoltaic system. To address communication challenges, the system employs a hybrid connectivity topology combining 4G networks and long-range radio communication (LoRa). The ecosystem is supported by dedicated software for the automated delivery of alerts and tactical dispatch of firefighting crews.
\end{abstract}
